\documentclass[11pt]{article}
\usepackage[margin=1in]{geometry}
\usepackage{amsmath,amssymb}
\usepackage[hidelinks]{hyperref}

\newcommand{\M}{\mathcal M}
\newcommand{\Hc}{\mathcal H}
\newcommand{\hc}{\mathsf h}
\newcommand{\Int}{\operatorname{Int}}

\title{Literature status: endpoint questions from arXiv:1712.01374}
\author{Run \texttt{fa\_banach\_001}}
\date{2026-08-09}

\begin{document}
\maketitle

\paragraph{Status.}
\texttt{literature\_already\_answered}.  This note records later literature
and makes no claim of a new proof.

\paragraph{Source questions.}
Randrianantoanina--Wu--Xu,
\emph{Noncommutative Davis type decompositions and applications},
arXiv:1712.01374, asks in Remark~3.11 whether its symmetric-space Davis
decomposition extends to every
\(E\in\Int(L_1,L_q)\), \(1<q<\infty\), and whether
\[
  \Hc_E^c(\M)=\hc_E^d(\M)+\hc_E^c(\M)
\]
holds for every \(E\in\Int(L_1,L_2)\).  The remark following its
Burkholder/Rosenthal theorem also asks whether either usual sum/intersection
formula applies to \(L_{2,\infty}\), or more generally \(L_{2,q}\) with
\(q\ne2\).

\paragraph{Explicit Davis answer.}
Randrianantoanina, \emph{P. Jones' interpolation theorem for noncommutative
martingale Hardy spaces}, arXiv:2212.08714, explicitly cites the source as
\texttt{Ran-Wu-Xu} and states that it resolves the problem left there.
Proposition~\texttt{Davis-symmetric}(i) proves, for every
\(E\in\Int(L_1,L_q)\),
\[
 E^{\rm ad}(\M;\ell_2^c)
 \subseteq
 E\!\left(\bigoplus_{n\ge1}\M_n\right)
 +E^{\rm cond,ad}(\M;\ell_2^c),
\]
which is the requested endpoint decomposition.  Part (ii) and
Corollary~\texttt{Davis} give the displayed Hardy-space identity for every
\(E\in\Int(L_1,L_2)\).  The same source proves this threshold sharp: validity
of the column Davis identity forces \(E\in\Int(L_1,L_2)\).

\paragraph{Weak-\(L_2\) answer.}
Cadilhac, \emph{Noncommutative Khintchine inequalities in interpolation
spaces of \(L_p\)-spaces}, arXiv:1809.06091, Proposition
\texttt{counterexample} and Section~6, constructs finite Schur--Horn examples
showing that neither standard row/column Khintchine formula holds in
\(L_{2,\infty}\).  Either Burkholder/Rosenthal alternative asked for in the
source would imply the corresponding Rademacher/Khintchine estimate; hence
the specifically highlighted weak-\(L_2\) case has a negative answer.
Cadilhac's arXiv:1812.03861, Theorem~\texttt{khrad}, gives the subsequent
strict-interpolation characterization of both standard Rademacher formulas.

\paragraph{Provenance.}
The 2022 paper knows that it is answering Remark~3.11 of the source.  The
weak-\(L_2\) conclusion is a direct consequence of the later explicit
Khintchine counterexamples and of the necessary reduction already noted in
the source paper.

\begin{thebibliography}{9}
\bibitem{RWX}
N. Randrianantoanina, L. Wu, and Q. Xu,
\emph{Noncommutative Davis type decompositions and applications},
arXiv:1712.01374.

\bibitem{Cad18}
L. Cadilhac,
\emph{Noncommutative Khintchine inequalities in interpolation spaces of
\(L_p\)-spaces}, arXiv:1809.06091.

\bibitem{CadMaj}
L. Cadilhac,
\emph{Majorization, interpolation and noncommutative Khintchine
inequalities}, arXiv:1812.03861.

\bibitem{JonesNC}
N. Randrianantoanina,
\emph{P. Jones' interpolation theorem for noncommutative martingale Hardy
spaces}, arXiv:2212.08714.
\end{thebibliography}

\end{document}

